\subsection*{Logarithmic derivatives}\label{sec:rem}
 The key fact about logarithmic derivatives that we are using is that if 
 $f(X), B(X) \in \F_p[X]$, $t \cdot \deg(f), \deg(B) < p$, and:
 $$ t \frac{f'(X)}{f(X)} = \frac{B'(X)}{B(X)},$$
 then $f(X)^t = \lambda \cdot B(X)$ for some constant $\lambda \in \F_p$. 
 
 We recap a quick proof. The hypothesis implies that $\left( \frac{f(X)^t}{B(X)} \right)' = 0$,
  and thus $\frac{f(X)^t}{B(X)}$ must be a rational function in $X^p$. To see the last deduction, note that $\left(f(X)^t \cdot B(X)^{p-1} \right)' = \left( \frac{f(X)^t}{B(X)} \cdot B(X)^p \right)' = \left( \frac{f(X)^t}{B(X)} \right)' \cdot B(X)^p + \frac{f(X)^t}{B(X)} \cdot p \cdot B(X)^{p-1} \cdot B'(X) = 0 + 0 = 0$. Thus the polynomial $f(X)^t \cdot B(X)^{p-1}$ is a polynomial in $X^p$, which implies that $\frac{f(X)^t}{B(X)} = 
 \frac{f(X)^t \cdot B(X)^{p-1}}{B(X)^p} =
 \frac{f(X)^t \cdot B(X)^{p-1}}{B(X^p)}$ is a rational function in $X^p$. 
 Once we know that $\frac{f(X)^t}{B(X)}$ is a rational function in $X^p$, our assumption about the degrees implies the result. 
